\begin{scholium}[7.8.3]
\EGAReaderTarget{ega4.c04.statement.7-8-3}{IV.2.s7.8.3-fr}
\textup{(i)} En las condiciones \textup{(i)} y \textup{(ii)} de la definición \textup{(7.8.2)}, basta considerar
los ideales $\mathfrak p$ maximales de $A$. Para que un anillo \textup{local}
noetheriano sea excelente, es necesario y
\oldpage[IV-2]{215}
suficiente que sea universalmente catenario y que sus fibras formales sean
geométricamente regulares.

\textup{(ii)} Si $A$ es un anillo excelente, también lo son toda localización
$S^{-1}A$ y toda $A$-álgebra de tipo finito.

\textup{(iii)} Un anillo local completo (en particular, un cuerpo) es excelente.
Un anillo de Dedekind cuyo cuerpo de fracciones tiene característica $0$
(en particular, $\mathbf Z$) es excelente.

\textup{(iv)} Sean $A$ un anillo excelente y $X=\operatorname{Spec}(A)$; el conjunto
$\operatorname{Reg}(X)$ \textup{(resp.} $\operatorname{Nor}(X)$,
$\mathrm U_{\mathrm R_n}(X)$) de puntos en los que $X$ es regular \textup{(resp.}
normal, \textup{resp.} satisface $(\mathrm R_n)$) es abierto en $X$. Para todo
$\mathcal O_X$-módulo coherente $\mathcal F$, el conjunto
$\mathrm U_{\mathrm C_n}(\mathcal F)$ \textup{(resp.}
$\mathrm U_{\mathrm S_n}(\mathcal F)$, $\operatorname{CM}(\mathcal F)$) de
los $x\in X$ en los que $\operatorname{coprof}(\mathcal F_x)\leq n$
\textup{(resp.} de los puntos en los que $\mathcal F$ satisface $(\mathrm S_n)$,
\textup{resp.} de los puntos en los que $\mathcal F_x$ es un $\mathcal O_x$-módulo
de Cohen-Macaulay) es abierto en $X$.

\textup{(v)} Sean $A$ un anillo excelente, $\mathfrak J$ un ideal de $A$ y
$\widehat A$ el completado separado de $A$ para la topología
$\mathfrak J$-preádica. Entonces el morfismo canónico
$f:\operatorname{Spec}(\widehat A)\to\operatorname{Spec}(A)$ es regular
\textup{(es decir,} plano \textup{y con fibras} geométricamente regulares). Si se pone
$X=\operatorname{Spec}(A)$, $X'=\operatorname{Spec}(\widehat A)$, se tiene
(con las notaciones de \textup{(iv)} y poniendo
$\mathcal F'=\mathcal F\otimes_{\mathcal O_X}\mathcal O_{X'}$):
\begin{env}[7.8.3.1]
\label{IV.2.s7.8.3.1-fr}
\begin{align*}
\operatorname{Reg}(X')&=f^{-1}(\operatorname{Reg}(X)),&
\operatorname{Nor}(X')&=f^{-1}(\operatorname{Nor}(X)),&
\mathrm U_{\mathrm R_n}(X')&=f^{-1}(\mathrm U_{\mathrm R_n}(X)),\\
\mathrm U_{\mathrm C_n}(\mathcal F')&=f^{-1}(\mathrm U_{\mathrm C_n}(\mathcal F)),&
\mathrm U_{\mathrm S_n}(\mathcal F')&=f^{-1}(\mathrm U_{\mathrm S_n}(\mathcal F)),&
\operatorname{CM}(\mathcal F')&=f^{-1}(\operatorname{CM}(\mathcal F))
\end{align*}
\end{env}

En particular, si $\mathfrak J$ está contenido en el radical de Jacobson de
$A$ \textup{(por ejemplo, si $A$ es local y $\mathfrak J$ es su ideal maximal)},
para que $A$ sea regular \textup{(resp.} normal, \textup{resp.} reducido, \textup{resp.} satisfaga
$(\mathrm R_n)$, \textup{resp.} tenga coprofundidad $\leq n$, \textup{resp.} satisfaga
$(\mathrm S_n)$, \textup{resp.} sea un anillo de Cohen-Macaulay), es necesario y
suficiente que $\widehat A$ tenga la misma propiedad; en particular, para
que $A$ no tenga ideales primos asociados inmersos, es necesario y suficiente
que ocurra lo mismo con $\widehat A$.

\textup{(vi)} Un anillo excelente $A$ es universalmente japonés; en particular, si
$A$ es un dominio de integridad, su cierre entero en cualquier extensión
finita de su cuerpo de fracciones es una $A$-álgebra finita.

\textup{(vii)} Sea $A$ un anillo \textup{local} excelente reducido. Entonces su completado
$\widehat A$ es reducido, el cierre entero $A'$ de $A$ en su anillo total
de fracciones es una $A$-álgebra finita \textup{(y, por tanto, un anillo} semilocal),
y el cierre entero de $\widehat A$ en su anillo total de fracciones es
isomorfo al completado $\widehat{A}'$ de $A'$. Además, existe una
correspondencia biyectiva canónica entre el conjunto de puntos maximales de
$\operatorname{Spec}(\widehat A)$ \textup{(es decir, el conjunto de ideales primos}
minimales \textup{de $\widehat A$)} y el conjunto de puntos cerrados de
$\operatorname{Spec}(A')$ \textup{(es decir, el conjunto de ideales} maximales \textup{de
$A'$)}. Para que el anillo local $A$ sea unirrama, es necesario y suficiente
que $\widehat A$ sea un dominio de integridad; para que $A$ sea
geométricamente unirrama, es necesario y suficiente que $\widehat A$ lo sea.

\textup{(viii)} Si $A$ es un anillo excelente que es un dominio de integridad, el
cierre entero de $A$ es la intersección de los cierres enteros de los anillos
$A_{\mathfrak p}$, donde $\mathfrak p$ recorre el conjunto de ideales
primos de $A$ de altura $1$.

\textup{(ix)} Sean $A$ un anillo excelente, $X=\operatorname{Spec}(A)$, $Z$ una parte
cerrada de $X$, $U=X\setmin Z$, $i:U\to X$ la inclusión canónica y $\mathcal F$
un $\mathcal O_U$-módulo coherente. Para que el $\mathcal O_X$-módulo
$i_*(\mathcal F)$ sea coherente, es necesario y suficiente que, para todo
$x\in\operatorname{Ass}(\mathcal F)$, se tenga
$\operatorname{codim}(\overline{\{x\}}\cap Z,\overline{\{x\}})\geq2$.
En particular, si $A$ es un dominio de integridad y $\mathcal F$ es no nulo y libre
de torsión,
para que $i_*(\mathcal F)$ sea coherente, es necesario y suficiente que
$\operatorname{codim}(Z,X)\geq2$.

\oldpage[IV-2]{216}

\textup{(x)} Sea $A$ un anillo local excelente. Para todo anillo cociente $B$ de $A$
que sea un dominio de integridad, $\widehat B$ es equidimensional. Para que
$A$ sea equidimensional, es necesario y suficiente que $\widehat A$ lo sea.
\end{scholium}

\begin{proof}
La mayoría de estos resultados ya se han demostrado.

(i) La primera afirmación se deduce de (5.6.3, (i)) y (7.4.4, (i));
la segunda, de (7.4.4), (7.3.18), (6.12.7) y (6.12.4).

La afirmación (ii) se deduce de (5.6.1), (5.6.3, (i)), (7.4.4) y (6.12.4).

(iii) La primera afirmación se ha visto en (5.6.4) y (7.4.4), teniendo en
cuenta (i). La segunda se ha visto en (5.6.4) y (7.3.19, (iv)), teniendo en
cuenta (i).

La afirmación (iv) se deduce de (6.12.4), (6.13.5), (6.12.9) y (6.11.8)
(teniendo en cuenta (ii)).

La primera afirmación de (v) es consecuencia de (7.4.6); de ella se deducen
las relaciones (7.8.3.1), teniendo en cuenta, respectivamente, (6.5.1),
(6.5.4), (6.5.3), (6.3.5) y (6.4.2). La última afirmación de (v) es el
caso particular correspondiente a la propiedad $(\mathrm S_n)$ para $n=1$
(5.7.5).

La afirmación (vi) se deduce de (7.7.3), y la afirmación (vii), de (7.6.1),
(7.6.2) y (7.6.3). Para demostrar (viii), observemos que el anillo
$A^{(1)}$, intersección de los $A_{\mathfrak p}$ para todos los ideales
primos $\mathfrak p$ de altura $1$ de $A$, es una $A$-álgebra \emph{finita}, como
se deduce de (vi), de (7.8.2, (i)) y de (5.11.2). Como el cierre entero
$A'$ de $A$ es una $A$-álgebra finita, los ideales primos de altura $1$ de
$A'$ son exactamente los que están por encima de los ideales primos de
altura $1$ de $A$ (5.10.17, (iv)); la conclusión se deduce, por tanto, de
(\textbf 0, 23.2.9).

La primera afirmación de (ix) es consecuencia de (vi) y de (5.11.4); la
segunda es un caso particular de ella, si se observa que, cuando $\mathcal F$
es no nulo y libre de torsión, $\operatorname{Ass}(\mathcal F)$ se reduce al punto
genérico de $\operatorname{Spec}(A)$. Por último, para demostrar (x), basta
observar que $B$ es un anillo local excelente por (ii); como $B$ es
universalmente catenario y universalmente japonés en virtud de (vi), el
anillo $B^{(1)}$ es una $B$-álgebra finita (5.11.2), lo que muestra que
$A$ es estrictamente formalmente catenario por (7.2.5, $c)$) y concluye la
demostración de (x).
\end{proof}
